\section{DiT、ConvNeXt 与反脆弱}

如果说前一章更多是在谈方法论，那么 ConvNeXt 与 DiT 则展示了这套方法论怎样在具体作品里落地。ConvNeXt 最有意思的地方，不是“卷积网络又赢了一次”，而是它敢于在 Vision Transformer 风头最盛的时候提出一个更尖锐的问题：ViT 到底为什么有效？谢赛宁与实习生刘壮合作时，并没有简单地在卷积网络上做怀旧式修补，而是反过来把 ViT 时代被证明有效的一系列训练和设计习惯重新移植回卷积框架中，再逐项分析哪些因素真正重要。最后得到的结论极具挑衅性：soft attention 也许不是 ViT 最关键的部分，很多收益来自更普适的训练配方、归一化、patchify 风格与 macro design。

\begin{knowledgebox}{ConvNeXt 的真正价值}
ConvNeXt 不只是“证明卷积还没死”，而是通过一组极有说服力的 ablation，拆开了 ViT 神话中的多个组成部分。它提醒研究者：当某个范式大获成功时，最重要的工作之一不是盲目模仿，而是辨认到底哪些成分在起作用，哪些只是时代风格的附属物。
\end{knowledgebox}

这种敢于挑战主流解释的姿态，在 DiT 上体现得更加明显。DiT 的起点并不是“我要发一篇生成模型爆款论文”，而是一个相当基础的问题：diffusion model 学到的表征，和自监督学习的表征相比到底怎样？在研究过程中，他们发现前者远不如后者，于是转而想：既然 diffusion pipeline 里的 U-Net 已经成为默认选项，为什么不能直接把 ViT 当作 denoiser backbone？结果非常简单，也非常有效。简单到什么程度？简单到论文投 CVPR 时被拒掉，理由之一是“太简单”。这一幕几乎像研究史对自身的讽刺：许多真正有力的想法，恰恰因为太直接、太少装饰，而不符合评审对“复杂创新”的想象。

\begin{importantbox}{DiT：简单结构的力量}
DiT 的核心并不在于引入了多少新模块，而在于它抓住了一个尺度化时代最关键的判断：如果 Transformer 在别的模态上已经证明了其统一性和扩展性，那么 diffusion 体系也应该允许更简洁、更可扩展的主干。越接近基础结构，越可能把系统推向更大的规模和更清晰的解释。
\end{importantbox}

后来 DiT 转投其他会议并成为 oral，LeCun 还专门发推吐槽 reviewer。谢赛宁借这段经历说出另一句很值得记住的话：\textit{``所有的 research paper 中或不中一点都不重要...完全是一个纯粹的随机过程''}。这里当然不是说评审毫无价值，而是说在前沿研究的边缘地带，很多工作之间的区分本来就极细，评审结果高度受具体 reviewer 口味、时间压力和背景知识影响。若把命中与否上升为对自我价值的判断，就很容易把研究节奏交给一个本质上高噪声的系统。更现实的是，DiT 诞生时 FAIR 已经开始文化转型，他离职后甚至不被允许署名，这又进一步说明，真正有影响的工作并不总能在当下获得“完美归属”。

\begin{warningbox}{把审稿结果当作真理，是研究者常见的自伤行为}
DiT 和 SiT 都经历了“先被拒、后被认可”的路径。谢赛宁想强调的不是自己受了委屈，而是研究者必须接受一个事实：评审制度只能粗糙抽样，无法完整测量工作价值。若把它误认为终审判决，就会把大量时间浪费在情绪性解释上，而不是继续向前推进。
\end{warningbox}

DiT 后来通过 Bill Peebles 进入 OpenAI 语境，并最终与 Sora 产生明显血缘，这让它更像一篇站在研究与工业边界上的论文。谢赛宁甚至把它只算作“0.25 篇”真正改变 AI 进程的 paper，意思并不是贬低自己，而是极其清醒地承认：很多工作是在大趋势已经形成时，把边界向前推了一小步。正因为有这种尺度感，他才能进一步提出“反脆弱”这个关键词。SiT 也是类似路径，flow matching 加 Transformer 的组合再次经历先拒后中，而这些随机冲击并没有摧毁研究轨迹，反而扩大了工作后来的影响。

\begin{importantbox}{Research 必须是反脆弱系统}
谢赛宁给出的定义非常实用：一个 random shock 带来的收益若大于损失，这个系统就是反脆弱的。对研究来说，这意味着你不能把全部价值压在某一次投稿、某个组织归属或某条单一路线上。相反，好的研究布局应当让失败、延迟、误判甚至拒稿，都有机会转化为更大的传播、重组或再发现。
\end{importantbox}

Perplexity 早期 demo 的故事也常被人拿来做 hindsight comedy：Aravind 在 Blue Bottle 咖啡店展示 demo 时，谢赛宁心想“这不就是 GPT 套个壳”，于是婉拒邀请。这个片段真正有意思的地方，不是证明谁看走了眼，而是说明未来往往由许多低维判断叠加而成。有人会错过显而易见的商业化机会，也有人会误判一项研究的长期价值；关键不在于永远押中，而在于让自己的能力结构不依赖每次都押中。反脆弱思维在这里再次成立：你必须构建一种即使偶尔看错，也不会被一次错失定义的人生与研究系统。

\subsection{本章小结}

ConvNeXt 和 DiT 共同说明，谢赛宁最擅长的不是追随热门，而是在热门范式内部追问“真正起作用的变量是什么”。配合对评审随机性的清醒认识和对反脆弱系统的强调，这一章展示的是一种能够承受冲击、并把冲击转化为增益的研究姿态。

